% !TeX root = ../papers/nc-p3.tex
\section{入力二十点からの全中央辺のGrassmann表示}
\label{sec:edge-plucker-presentations}
\begin{theorem}\label{thm:edge-plucker-presentations}
有限辺の一位置の点窓 $X\subset\mathbb P^3$ を固定し、$W=H^0(\mathcal I_X(4))$ とする。
$X$ を入力に持つ有限辺全体は、六次元空間 $R$、$\bigwedge^2 R$ に値を持つ四次式列 $P(x)$ の
次の条件によって、出力座標変更を除き過不足なく復元できる。
\begin{enumerate}
\item 十五成分は $W$ の基底で $P\wedge P=0$ である。
\item $r\in R^*\otimes S_1$ について $\iota_r P=0$ を多項式恒等式とする線型解の中に、
分数体上で独立な四行がある。
\end{enumerate}
同一の辺を与える二表示は $\mathrm{GL}(R)$ の作用で移り合う。
線型解は条件成立時ちょうど四次元で、その基底が辺の一次四行六列行列になる。
入力の標識も忘れれば、更に $X$ の射影自己同型の作用で割る。
\end{theorem}
\begin{proof}
有限辺の六つの双一次関係を $R$ とし、入力 $x$ で評価した行列を $A(x):R\to k^4$ とする。
長さ20の期待余次元なのでEagon--Northcott分解により十五の最大行列式が四次の基底を与え、
イデアルはこれらで生成される。符号付き相補列で並べれば $P$ は一般核のPlücker二ベクトルであり、
四行は全てこれを縮約して零になる。

逆に独立四行を取ると、分数体上その共通二次元核のPlücker座標は $P$ に比例する。
元の十五成分は有限基点しか持たず原始であり、四行の最大行列式も次数4なので比例係数は非零定数である。
従って同じ飽和イデアル $I_X$ を得る。この四行以外の線型解は、$X$ の外で全階数の四行の正則関数係数の
線型結合である。両辺の次数が1で係数は次数零、余次元3の補集合から延びるので全て定数である。
また六関係は独立である。一定の列核があれば $P$ の値は固定ベクトルとの外積に入り、十五成分独立に反する。

$X$ の任意の台点で行列の階数は3である。実際階数が $s\le2$ なら単元 $s$ 行列を消去した後の
$(4-s)\times(6-s)$ の最大行列式イデアルは局所環上で期待余次元3を持つ。
残る全成分は極大イデアルに入り局所Eagon--Northcott分解が極小だから、最小生成元は
$\binom{6-s}{4-s}>3$ 個で、局所完全交叉 $X$ に反する。従って入力射影での核直線（四行の転置の核）は
$X$ 上の正則射を与え、六双一次式の零スキームはそのグラフである。他方の射影でも正次元の
線型繊維は存在せず、同じ単元小行列式による消去で閉埋込みとなる。
二つの関係基底の違いがちょうど $\mathrm{GL}(R)$、四行の基底の違いが出力の座標変更である。
この結果は辺の完全な認識であり、正則な四窓への延長には切断の条件がさらに必要となる。
\end{proof}

\begin{corollary}\label{cor:general-edge-two-grassmannians}
十五次元成分の一般窓の一つの入力配置で、条件を満たすGrassmann多様体はちょうど二つである。
それらは二番目の線型解の条件を外した全Grassmann包含の中でも、被約な孤立点である。
\end{corollary}
\begin{proof}
任意の表示の出力 $Y$ について、有限辺の完全交叉の跡により $Y^2$ は入力の平方空間の直交補に
単元を掛けたものとなる。完全系列の二次評価が十次元であることも、両辺の最大行列式分解から従う。
既知の順・逆辺の二出力は射影非同値で、双方の間の反対行列の階数2以下軌跡は有限である。
従って定理\ref{thm:clean-configuration-extensions}の平方根の一意性の議論で、
全階数の射影平方根はその二つだけである。入力と出力の二射が決まれば
グラフ上で消える $(1,1)$ の六式がちょうど関係であるから表示もこの二つに限られる。

第一条件の可逆な十五座標の変形を書き、一次項を $D\in\operatorname{End}(\bigwedge^2 R)$ とする。
場合分けのない接条件は $P(x)\wedge D P(x)=0$ である。$\mathfrak{gl}_6$ による36方向は常に入り、
十五成分が独立なのでこの作用の微分は忠実である。
\path{computations/verify_edge_plucker_presentations.py} は既知の有理標本の良い素数2657で、
四辺につき百個の入力評価だけで一次条件の階数189を確認する。従って全係数での核はちょうど36次元である。
これは開な小行列条件で、標数零の一般窓でも成立する。
$\mathrm{GL}_{15}$ の内部のこの座標表示から同一Grassmann多様体の表示への繊維は
$\bigwedge^2\mathrm{GL}_6$ である。実際 $\operatorname{Gr}(2,6)$ の射影自己同型は
二種の極大線型空間の次元が異なるので基礎射影五空間の射影自己同型に限られる。
滑らかなこの群の商の接空間が零となり、主張を得る。
\end{proof}

同じ検算で四辺の $\iota_rP$ の評価階数は20となり、二十四未知数から元の四行が回収される。
また四次基底の積には85独立なものがある。点窓の二重化の分解から八次部分は
$10\binom{5}{3}-15=85$ 次元なので、この標本と一般点では標準四次埋込みの二次イデアルは三十五次元である。
